\section{Metrics: Proposal and Critical Evaluation (RQ2, Part 2)}
\label{sec:metrics}

We evaluate six candidate metrics---and commit in advance to rejecting any that
cannot be justified---against three criteria: \emph{grounding}
(does it estimate a quantity of \cref{sec:definitions,sec:model}?),
\emph{operationalizability} (can it be computed from obtainable data?), and
\emph{incentive safety} (what does optimizing it do?). One candidate is rejected;
one is accepted only with a replacement semantics; the remainder are adopted with
stated scope. The foundational metrics (EDR analogues, ERC/FRR/AR reconstruction performance,
irrecoverability, and $\ED$ itself) are computed concretely in
\cref{sec:simulation}; the replacement decay-schedule semantics is computed in
\cref{sec:interest}; the flow metric (M4) and the composite (M6) are defined for
field use and not exercised here.

\begin{metricdef}[Evidence Deficit Ratio, EDR]
\label{met:edr}
For an estate and schema $\Sch$: the fraction of schema-required elements and link
keys absent from the corpus, reported per deficit class (missing node, missing key,
degraded attribute) and per interrogative (who/what/why/authorization/outcome/%
recovery). \emph{Verdict: adopted.} Grounded directly in \cref{def:deficit};
computable by automated census; class-resolved reporting resists the false comfort
of a single number. This is the suite's foundation metric---syntactic, cheap, and
continuously monitorable.
\end{metricdef}

\begin{metricdef}[Evidence Reconstruction Performance, ERC, FRR, and AR]
\label{met:erc}
The triplet comprising (i)~Evidence Reconstruction Cost (ERC): measured person-time,
compute, or calendar latency for a defined query; and (ii)~False Reconstruction
Rate (FRR): the fraction of accepted facts or links that diverge from independent
ground truth; and (iii)~policy Abstention Rate (AR): the fraction for which $\rho$
certifies no answer. All are measured under the same preregistered confidence and
acceptance policy, per \cref{def:rec}, and collected from
real reconstructions (incidents, audits) or standardized drills (\cref{sec:field});
in simulation, effort is proxied by candidate comparisons and FRR is scored exactly.
\emph{Verdict: adopted.} The triplet exposes the effort--accuracy--abstention
tradeoff that ERC alone can hide. AR is not IRR: the former is a policy outcome,
the latter a certified evidence state.
\end{metricdef}

\begin{metricdef}[Irrecoverability Ratio, IRR]
\label{met:irr}
The fraction of a defined query set (or of required links) irrecoverable at $t$
within a published $(U,\gamma,\rho_{\mathrm{exh}})$ declaration, per
\cref{def:irrec}. \emph{Verdict: adopted.} Decidable exactly in simulation;
certifiable in the field by source-exhaustion protocols~\cite{nist80086}. Reported
separately from cost metrics because it measures loss, not effort, and the two must
never be averaged together (\cref{prop:writeoff}).
\end{metricdef}

\begin{metricdef}[Evidence Debt Inflow Rate, EDIR]
\label{met:growth}
The inflow attributable to \emph{new} deficit creation per unit time:
in practice, the deficit-census delta per unit time (new changes arriving with
incomplete evidence), weighted by principal estimates. \emph{Verdict: adopted.}
This measures the inflow (are we borrowing faster?), and must
be reported separately from interest accrual on the existing stock; conflating the
two makes remediation invisible. The net stock change is instead
$d\ED/dt=\text{inflow}+\text{decay accrual}-\text{repayment}
-\text{write-off realization}$.
\end{metricdef}

\begin{metricdef}[Evidence-Debt Interest Rate]
\label{met:interest}
\emph{Proposed as}: a scalar rate $\rho$ such that reconstruction cost grows as
$e^{\rho\tau}$ or linearly at rate $\rho$. \emph{Verdict: rejected.}
\Cref{prop:steps} shows that growth in effort and abstention exposure under sound
exhaustive policies is schedule-driven and generally
piecewise, while error depends on discriminator degradation; neither supports a
single invariant rate. Any fitted scalar is an artifact of the observation window and
extrapolates falsely across retention cliffs. \emph{Replacement:} the
\textbf{decay schedule} of the estate---the dated sequence of upcoming path-death
events (retention expirations, migrations, projected turnover quantiles) with the
cost jump each implies---and, where a single summary is required, the
\textbf{evidence half-life}: the age $\tau_{1/2}$ at which half of a reference
query set ceases to be answerable at baseline cost. Both are computable from the
survival inventory (\cref{prop:reduction}) and are demonstrated in
\cref{sec:interest}.
\end{metricdef}

\begin{metricdef}[Evidence Debt Index, EDI]
\label{met:edi}
A composite scalar summarizing estate debt for management reporting: a weighted
aggregation of EDR (by interrogative), calibrated ERC, false-reconstruction
exposure (FRR priced by $\lambda_q$), abstention exposure (AR priced by $\pi_q$),
projected decay from the decay schedule, and separately reported IRR, with published
weights. \emph{Verdict: adopted with explicit estimator
status and health warnings.} By \cref{sec:model-notformula}, any such composite is
a policy-specific estimator related to $\ED_\rho$ under assumptions, not a
definition of $\ED^{\star}$;
its weights encode the workload judgment $(w_q, \lambda_q, \pi_q)$ and must be published for
any comparison; and it inherits the gaming risks of all composites---optimizing EDI
by weakening the schema $\Sch$ is the exact analogue of declaring bankruptcy
solvent by redefining assets. EDI is defensible only alongside the disaggregated
vector (EDR by class, measured ERC/FRR/AR, decay schedule, and IRR).
\end{metricdef}

Two measurement caveats apply to the whole suite. \emph{Discoverability bias}: the
census can only inspect stores it knows about; evidence in unknown stores is
neither counted as present nor as deficit---field deployments should treat store
discovery as part of the census. \emph{Confidence accounting}: reconstructed
answers vary in confidence; ERC and FRR must be reported under the same stated
confidence and acceptance policy, or performance is not comparable. The simulation sidesteps both (closed
world, exact scoring); the field protocol confronts them directly
(\cref{sec:field}).
